% ===== 原文 PDF pp.41--43（印刷 pp.40--42） =====
% tr-04 的 tr04a 片段。起点：tr-03 已写完 p.40（印刷 p.39）上开始的句子，
% p.41 首行即「例 3.1」，无跨页半句需要续写。
% 本文件 tag：3.28--3.31（印刷 p.40）；3.32--3.35（印刷 p.42）。共 8 条。

% [原文 p.40]
\begin{examplebox}{例 3.1：}
我们给出一个例子，说明单值性如何刻画库仑分支上的奇点。考虑一个 $\mathrm{U}(1)$ 矢量多重态与一个电荷为 $\sqrt{p}$ 的超多重态耦合的情形。单圈跑动给出
\begin{equation}
  \tau(a) = \frac{p}{2\pi \ii}\log\left(\frac{a}{\Lambda_{\mathrm{UV}}}\right),
  \tag{3.28}
\end{equation}
其中 $a$ 是矢量多重态的标量，它给出带电超多重态的质量。特别地，若我们在库仑分支上移动到 $a\sim 0$ 附近——此处超多重态变为无质量——并在绕 $a=0$ 转一圈之后，就会看到 (3.28) 中的耦合获得一个平移
\begin{equation}
  \tau \to \tau + p \,.
  \tag{3.29}
\end{equation}
而由式 (3.20) 我们知道，相应的 $\SL(2,\Z)$ 单值性为
\begin{equation}
  M = \begin{pmatrix}
    1 & p \\
    0 & 1
  \end{pmatrix}.
  \tag{3.30}
\end{equation}
\end{examplebox}

% [Tutor Note: (3.28) 中的 $p$ 是超多重态电荷的平方：电荷写成 $\sqrt{p}$，单圈跑动的系数正是 $p/(2\pi\ii)$。绕 $a=0$ 一周时 $\log a$ 多出 $2\pi\ii$，于是 $\tau\to\tau+p$，这就是 (3.29)。]
% [Tutor Note: (3.30) 的矩阵作用在列向量 $(a_D,a)^{\mathsf T}$ 上（$a_D$ 在上，与式 (3.26) 的写法一致）：这时 $a_D\to a_D+p\,a$、$a\to a$，于是 $\tau=\partial a_D/\partial a\to\tau+p$，与 (3.29) 自洽。若误以为它作用在 $(a,a_D)^{\mathsf T}$ 上，会得到 $\tau\to\tau/(1+p\tau)$，与 (3.29) 矛盾。]

更一般地，还可以有别的粒子变成无质量的奇点：不只是电子，更一般地还有带电磁荷 $(p,q)$ 的双子。它们的每一个单值性都与超多重态中电子的单值性属于同一共轭类，只不过由不同的 $\SL(2,\Z)$ 元素表示。相应的单值性为
\begin{equation}
  M^{p,q} = \begin{pmatrix}
    1+pq & p^2 \\
    -q^2 & 1-pq
  \end{pmatrix}.
  \tag{3.31}
\end{equation}
另一种可能是：多个 BPS 粒子在奇点处同时变成无质量，这会产生更一般的单值性，它们由 Kodaira 分类来刻画。更多信息可参阅 [163]。其中对我们后文讨论重要的是这样一类情形：两个 BPS 粒子的电荷分别为 $(p,q)$ 与 $(p',q')$，它们是相互非局域的，即 $(pq'-p'q)\neq 0$。这样的奇点给出一个强耦合的 SCFT，即所谓的 Argyres--Douglas 理论，我们稍后还会回到它。

% [原文 p.41]
\begin{center}
  \includegraphics[width=0.6\linewidth]{fig6.png}
\end{center}
\figtitle{6}{$u$ 平面上涌现的椭圆纤维化。}

% [Tutor Note: 图 6 中，$\cM_1$、$\cM_2$、$\cM_\infty$ 三点位于 $u$ 平面（右下角标注 $(a,a_D)$）上，虚线表示各自“长出”的椭圆纤维 $\Sigma_u$；蓝圈内为强耦合区，圈内叉号标出奇点，$\cM_\infty$ 位于弱耦合端。图内英文标注为原图固有，不重绘、不擦除。]

\SUB{3.3}{量子预势的 Seiberg--Witten 解}

电磁单值性告诉我们，特殊坐标 $(a,a_D)$ 在 $\cM_{\mathrm{CB}}$ 中某个奇点附近是如何变化的。然而，要确定 $F_k$、从而确定预势 $\mathcal{F}$，就必须知道 $(a,a_D)$ 的精确表达式。现在的问题是：要找出这样的多值函数 $(a,a_D)$，使它在 $\cM_{\mathrm{CB}}$ 的每个奇点附近都呈现出所需的单值性 $M$。这被称为 Riemann--Hilbert 问题；在乘以一个整函数的意义下，它的解是唯一的。事实证明，这样的解 $(a,a_D)$ 具有很好的几何图像。

主要的出发点是：有效理论的 $\SL(2,\Z)$ 不变性促使 Seiberg 与 Witten 引入一个额外的几何对象——环面 $T^2$，$\SL(2,\Z)$ 在其上作为模对称群作用。\footnote{这里我们关注的是秩一情形。对于一般的秩 $r$ 规范群，环面将被一个亏格为 $r$ 的黎曼面取代，见第 4 节。}更确切地说，考察 $\cN=2$ 理论的库仑分支（如图 6 所示意），很自然地可以把库仑分支上的有效耦合 $\tau^{\mathrm{IR}}$ 解释为一个“涌现”环面的复结构。带有全纯截面的代数环面——通常称为椭圆曲线，记作 $\Sigma_u$——可以代数地描述为

% [原文 p.42]
Weierstrass 模型，它由 $\mathbb{P}^2$ 中的一张超曲面定义：
\begin{equation}
  y^2 = x^3 + f(u,m)x + g(u,m)\,.
  \tag{3.32}
\end{equation}
复结构 $\tau^{\mathrm{IR}}$ 依赖于 $u$——$\cM_{\mathrm{CB}}$ 上的整体坐标——而这种依赖被编码在式 (3.32) 的系数 $f$ 与 $g$ 之中。参数 $m$ 表示其他可被开启的质量形变。这样的曲线 $\Sigma_u$ 就称为 Seiberg--Witten 曲线。

就这条 Seiberg--Witten 曲线而言，其上存在由 1-闭链组成的典范基 $\{A,B\}\in H_1(\Sigma_u,\Z)$；我们可以把库仑分支上的特殊坐标 (3.23) 认定为周期积分
\begin{equation}
  a = \oint_A \lambda\,,\qquad a_D = \oint_B \lambda\,,
  \tag{3.33}
\end{equation}
其中 $\lambda$ 是某个 1-形式微分，称为 \textbf{SW 微分}。它是 (3.32) 上的亚纯 1-形式，并服从特殊 Kähler 约束，即
\begin{equation}
  \frac{\partial \lambda}{\partial u} = \Omega + d\phi = \frac{dx}{y} + d\phi\,,
  \tag{3.34}
\end{equation}
其中 $\Omega=\frac{dx}{y}$ 是 SW 曲线上全纯且处处非零的 1-形式，$d\phi$ 指任意的恰当 1-形式。此外，$\lambda$ 的留数依赖于 EFT 中可能的质量形变。规范耦合 $\tau^{\mathrm{IR}}$ 等同于 $\Sigma_u$ 的复结构，并由如下周期之比给出
\begin{equation}
  \tau^{\mathrm{IR}} := \frac{\partial a_D}{\partial a}
  = \dfrac{\oint_A \Omega}{\oint_B \Omega}\,.
  \tag{3.35}
\end{equation}

% [Tutor Note: (3.34) 就是 SW 微分的“特殊 Kähler 约束”：$\partial_u\lambda$ 与全纯 1-形式 $\Omega=dx/y$ 只相差一个恰当形式 $d\phi$。恰当形式沿闭链积分给零，故它不改变 $a$、$a_D$（从而不改变 $\tau^{\mathrm{IR}}$），只影响 $\lambda$ 的留数，即质量项。]
% [Tutor Note: (3.35) 的分子/分母与 (3.33) 的定义互为倒数：由 $a=\oint_A\lambda$、$a_D=\oint_B\lambda$ 得 $\partial a_D/\partial a=(\oint_B\partial_u\lambda)/(\oint_A\partial_u\lambda)=\oint_B\Omega/\oint_A\Omega$，而原文写的是 $\oint_A\Omega/\oint_B\Omega$（已 700 dpi 裁图复核，底本确实如此）。故此处疑为原文笔误；译文照录原式，读者请以 (3.33) 的定义为准。若把基换成 $(-B,A)$（保持交截数 $A\cdot B=1$），则 $\tau^{\mathrm{IR}}\to-1/\tau^{\mathrm{IR}}$，正是电磁对偶的 S 变换。]

BPS 粒子的质量由 $\lambda$ 在 $\Sigma_u$ 上某条非平凡 1-闭链上的积分给出。我们在前几节已经看到，库仑分支上的奇点可以解释为某些 BPS 粒子变为无质量的点。就 $\Sigma_u$ 上的闭链而言，这意味着奇点就是库仑分支上相应闭链退化的点。而与这种奇点相关的单值性有一个很好的几何解释，即 Picard--Lefshetz 变换。也就是说，在一个 $A$-闭链消失的奇点附近，
